\documentclass[a4paper,12pt]{article}

\usepackage[utf8]{inputenc}
\usepackage[T1]{fontenc}
\usepackage[italian]{babel}
\usepackage{geometry}
\usepackage{amsmath}
\usepackage{booktabs}
\usepackage{xcolor}
\usepackage{listings}
\usepackage{enumitem}

\geometry{margin=2.5cm}

\lstset{
    basicstyle=\ttfamily\small,
    backgroundcolor=\color{gray!10},
    frame=single,
    breaklines=true,
    columns=fullflexible
}

\title{Interpretazione del Decay Factor} 
\author{}

\begin{document}

\maketitle
\newpage

\section{Interpretazione del Decay Factor}
Fino ad ora si è considerata la Tabù Matrix come un modo per penalizzare tutte le soluzioni già viste o scartate. In questo modo
le penalizzazioni vecchie restano forti anche quando magari non contano più tanto. 
Una variazione proposta potrebbe essere quella di introdurre un Decay Factor: ovvero un fattore di decadimento che permette di cacellare 
la memoria passata ad ogni iterazione, in particolare si è suggerito di adottare:
\[
    DF = 0.97
\] 
in questo modo permettiamo alla Tabù Matrix di dimenticare ad ogni iterazione il 3\% della memoria passata, facendo pesare meno le vecchie penalizzazioni
e di più quelle attuali.

\end{document}